%% kapa: exam=1 round=36 year=2025 period=1 form=A subject=civil_law subject_ko=민법 seq=29 number=29
%% source: Q-Net 감정평가사 기출문제 (cst00309) artlSeq=5231525 file=raw/1008/5231525/2025년 제36회 감정평가사 1차시험 1교시 문제지 원본.pdf
%% conversion: kordoc 4.24.51 -> tools/gen1.py (auto); review=proofread-sonnet; flags=none
\begin{question}{29}
취득시효에 관한 설명으로 옳지 않은 것은? (다툼이 있으면 판례에 따름)
\begin{choices}
  \item 국유재산 중 행정재산은 공용이 폐지되지 않는 한 취득시효의 대상이 되지 않는다.
  \item 성명불상자의 소유물에 대하여도 시효취득을 인정할 수 있다.
  \item 부동산에 관하여 적법한 등기를 마치고 소유권을 취득한 자가 그 부동산을 점유하는 경우, 특별한 사정이 없는 한 그 점유는 취득시효의 기초가 되는 점유라 할 수 없다.
  \item 1필의 토지의 일부분에 대한 시효취득도 인정될 수 있다.
  \item 1동의 집합건물의 구분소유자들의 건물 대지 전체에 대한 공동의 점유는 대지 소유권의 점유취득시효의 요건인 점유에 해당하지 않는다.
\end{choices}
\end{question}
\answer{5}
